
Large language models trained on web-scale corpora may overlap with standardized evaluation benchmarks, potentially inflating reported performance [Brown et al., 2020; Shi et al., 2024; Deng et al., 2024; Yang et al., 2023]. Multiple detection paradigms have been proposed to identify such contamination: n-gram overlap [Brown et al., 2020; Gao et al., 2021], embedding similarity [Yang et al., 2023; Singh et al., 2024], and membership inference attacks (MIAs) such as Min-K\% Prob \citep{Shietal2024}, Min-K\%++ [Zhang et al., 2024a], and DC-PDD [Zhang et al., 2024b]. These methods have been observed to produce inconsistent signals on the same benchmarks [Singh et al., 2024; Fu et al., 2024], without a principled framework explaining when each method is expected to succeed or fail.

The three distinct signal types operated upon by these detector families — lexical overlap, semantic proximity, and memorization signatures through model log-probabilities — are in principle orthogonal. A benchmark item's position in the two-dimensional space of (lexical overlap, semantic proximity) could, in principle, predict which detector family will perform best on that item. This motivates a three-zone phase diagram for contamination detection routing.

This paper describes hypothesis H-GeomRoute-v1, which proposes that corpus-side geometric signals (max 13-gram overlap count and SBERT cosine similarity to the nearest corpus neighbor) can classify benchmark items into contamination geometry strata, and that a logistic regression routing rule trained on one corpus will predict the top-performing detector family with cross-corpus top-1 accuracy > 40\% and Kendall's τ above a simulation-calibrated threshold on determinate items.

The executed experiment (h-e1, EXISTENCE sub-hypothesis) does not produce routing accuracy measurements. Instead, the experiment encounters a structural failure mode: when SBERT cosine similarity is computed against randomly-streamed corpus documents rather than top-k nearest-neighbor retrieved documents, the resulting similarity distribution is degenerate (near-zero) for all benchmark items at 50,000-document scale, causing all 25,403 items to collapse into the lexical stratum. This failure mode is termed stratum collapse.

The paper documents the stratum collapse finding, its root cause, the benchmark-corpus overlap characterizations produced as a byproduct, and four implementation bugs that compound the detector results. The main hypothesis remains unresolved: the gate failure is attributed to experimental design and implementation failures, not to falsifying evidence against the hypothesis.

The paper is organized as follows. Section 2 surveys related work. Section 3 describes the three-zone phase diagram methodology. Section 4 details the experimental setup. Section 5 presents results. Section 6 discusses implications and limitations. Section 7 concludes.

